\documentclass{article}
\usepackage[T1]{fontenc}
\usepackage{lmodern}
\usepackage{graphicx}
\usepackage{amsmath,amssymb}
\usepackage[a4paper,margin=2cm]{geometry}
\usepackage[absolute,overlay]{textpos}
\setlength{\TPHorizModule}{1mm}
\setlength{\TPVertModule}{1mm}
\usepackage[indonesian]{babel}
\usepackage{hyperref}
\setlength{\emergencystretch}{2em}
% unit: Halaman 52

\begin{document}

% === Halaman 52 (python/native) ===

52

Serangan lainnya: Side-Channel Attack

• Serangan ini tidak terutama menyerang algoritma matematisnya, tetapi memanfaatkan 

informasi yang bocor ketika algoritma dijalankan pada perangkat fisik.

• Ada beberapa macam side-channel attack:

1. Timing Attack

Memanfaatkan perbedaan waktu eksekusi untuk memperoleh informasi tentang kunci. 

2. Power Analysis Attack

Menganalisis konsumsi daya perangkat ketika melakukan operasi kriptografi. Contohnya: 

o

Simple Power Analysis (SPA) 

o

Differential Power Analysis (DPA) 

3. Electromagnetic (EM) Attack

Menganalisis radiasi elektromagnetik yang dihasilkan perangkat ketika melakukan operasi 

kriptografi. 

4. Fault Injection Attack

Penyerang sengaja menyebabkan kesalahan pada perangkat atau proses komputasi, 

kemudian menganalisis output yang salah tersebut. Salah satu contohnya adalah 

Differential Fault Analysis (DFA).

\end{document}
